\documentclass[12pt,a4paper]{article}
\usepackage{amsmath,amsfonts,amssymb,mathrsfs,tikz,times,pifont}
\usepackage{enumitem}
\newcommand\circitem[1]{%
\tikz[baseline=(char.base)]{
\node[circle,draw=gray, fill=red!55,
minimum size=1.2em,inner sep=0] (char) {#1};}}
\newcommand\boxitem[1]{%
\tikz[baseline=(char.base)]{
\node[fill=cyan,
minimum size=1.2em,inner sep=0] (char) {#1};}}
\setlist[enumerate,1]{label=\protect\circitem{\arabic*}}
\setlist[enumerate,2]{label=\protect\boxitem{\alph*}}
%%%::::::by chnini ameur :::::::%%%
\everymath{\displaystyle}
\usepackage[left=1cm,right=1cm,top=1cm,bottom=1.7cm]{geometry}
\usepackage{array,multirow}
\usepackage[most]{tcolorbox}
\usepackage{varwidth}
\usepackage{mathrsfs} % pour \mathscr
\usepackage{calligra}
\newcommand{\Dcalligra}{\text{\calligra D}}
\tcbuselibrary{skins,hooks}
\usetikzlibrary{patterns}
%%%::::::by chnini ameur :::::::%%%
\newtcolorbox{exa}[2][]{enhanced,breakable,before skip=2mm,after skip=5mm,
colback=yellow!20!white,colframe=black!20!blue,boxrule=0.5mm,
attach boxed title to top left ={xshift=0.6cm,yshift*=1mm-\tcboxedtitleheight},
fonttitle=\bfseries,
title={#2},#1,
% varwidth boxed title*=-3cm,
boxed title style={frame code={
\path[fill=tcbcolback!30!black]
([yshift=-1mm,xshift=-1mm]frame.north west)
arc[start angle=0,end angle=180,radius=1mm]
([yshift=-1mm,xshift=1mm]frame.north east)
arc[start angle=180,end angle=0,radius=1mm];
\path[left color=tcbcolback!60!black,right color = tcbcolback!60!black,
middle color = tcbcolback!80!black]
([xshift=-2mm]frame.north west) -- ([xshift=2mm]frame.north east)
[rounded corners=1mm]-- ([xshift=1mm,yshift=-1mm]frame.north east)
-- (frame.south east) -- (frame.south west)
-- ([xshift=-1mm,yshift=-1mm]frame.north west)
[sharp corners]-- cycle;
},interior engine=empty,
},interior style={top color=yellow!5}}
%%%%%%%%%%%%%%%%%%%%%%%

\usepackage{fancyhdr}
\usepackage{eso-pic}         % Pour ajouter des éléments en arrière-plan
% Commande pour ajouter du texte en arrière-plan
%\AddToShipoutPicture{
%    \AtTextCenter{%
%        \makebox[0pt]{\rotatebox{80}{\textcolor[gray]{0.7}{\fontsize{5cm}{5cm}\selectfont PGB}}}
%    }
%}
\usepackage{lastpage}
\fancyhf{}
\pagestyle{fancy}
\renewcommand{\footrulewidth}{1pt}
\renewcommand{\headrulewidth}{0pt}
\renewcommand{\footruleskip}{10pt}
\fancyfoot[R]{
\color{blue}\ding{45}\ \textbf{2025}
}
\fancyfoot[L]{
\color{blue}\ding{45}\ \textbf{Dindéfélo-Ségou}
}
\cfoot{\bf
\thepage /
\pageref{LastPage}}
\begin{document}
\renewcommand{\arraystretch}{1.5}
\renewcommand{\arrayrulewidth}{1.2pt}
\begin{tikzpicture}[overlay,remember picture]
\node[draw=blue,line width=1.2pt,fill=purple,text=blue,inner sep=3mm,rounded corners,pattern=dots]at ([yshift=-2.5cm]current page.north) {\begingroup\setlength{\fboxsep}{0pt}\colorbox{white}{\begin{tabular}{|*1{>{\centering \arraybackslash}p{0.28\textwidth}} |*2{>{\centering \arraybackslash}p{0.2\textwidth}|} *1{>{\centering \arraybackslash}p{0.19\textwidth}|} }
\hline
\multicolumn{3}{|c|}{$\diamond$$\diamond$$\diamond$\ \textbf{Dindéfélo-Ségou}\ $\diamond$$\diamond$$\diamond$ }& \textbf{A.S. : 2025/2026} \\ \hline
\textbf{Matière: Mathématiques}& \textbf{Niveau : 4}\textbf{ème} &\textbf{Date: 17/12/2025} & \textbf{Durée : 2 heures} \\ \hline
\multicolumn{4}{|c|}{\parbox[c]{10cm}{\begin{center}
\textbf{{\Large\sffamily Composition Du 1$ ^\text{\bf er} $ Semestre}}
\end{center}}} \\ \hline
\end{tabular}}\endgroup};
\end{tikzpicture}
\vspace{3cm}
\section*{Exercice 1 : Identités remarquables et Factorisation (5 points)}
\begin{enumerate}
    \item En utilisant les identités remarquables $(a+b)^2$, $(a-b)^2$ ou $(a-b)(a+b)$, développer et réduire les expressions suivantes : \textbf{(2 pts)}
    \begin{enumerate}
        \item $A = (x + 5)^2$
        \item $B = (3x - 2)^2$
        \item $C = (x - 4)(x + 4)$
    \end{enumerate}
    
    \item Développer et réduire les expressions suivantes : \textbf{(2,5 pts)}
    \begin{enumerate}
        \item $A = 3(2x - 5) + 4x(x + 2)$
        \item $B = (2x + 3)(x - 4) - (x^2 - 5)$
    \end{enumerate}

    \item Factoriser les expressions suivantes en recherchant un facteur commun : \textbf{(2,5 pts)}
    \begin{enumerate}
        \item $E = 5x^2 - 15x$
        \item $F = (x + 3)(2x - 1) + (x + 3)(x + 5)$
        \item $G = (x - 2)^2 - 3(x - 2)$
    \end{enumerate}
\end{enumerate}
\section*{\underline{Exercice 1 :} Développement et Factorisation (5 points)}
\begin{enumerate}
    \item Développer et réduire les expressions suivantes : \textbf{(2,5 pts)}
    \begin{enumerate}
        \item $A = 3(2x - 5) + 4x(x + 2)$
        \item $B = (2x + 3)(x - 4) - (x^2 - 5)$
    \end{enumerate}
    \item Factoriser les expressions suivantes en recherchant un facteur commun : \textbf{(2,5 pts)}
    \begin{enumerate}
        \item $C = 5x^2 - 15x$
        \item $D = (x + 3)(2x - 1) + (x + 3)(x + 5)$
        \item $E = 7x(x - 2) - (x - 2)$
    \end{enumerate}
\end{enumerate}

\section*{\underline{Exercice 2 :} Équations et Inéquations (5 points)}
\begin{enumerate}
    \item Résoudre dans $\mathbb{R}$ les équations suivantes : \textbf{(2,5 pts)}
    \begin{enumerate}
        \item $5x - 8 = 2x + 7$
        \item $3(x - 2) = 4x + 5$
        \item $(2x - 4)(3x + 9) = 0$
    \end{enumerate}
    \item Résoudre dans $\mathbb{R}$ les inéquations suivantes et représenter les solutions sur une droite graduée : \textbf{(2,5 pts)}
    \begin{enumerate}
        \item $2x + 5 < 13$
        \item $-3x + 4 \geq 10$
    \end{enumerate}
\end{enumerate}

\section*{\underline{Exercice 3 :} Droites remarquables du triangle (5 points)}
\begin{enumerate}
    \item \textbf{Définitions :} Répondre par vrai ou faux. \textbf{(2 pts)}
    \begin{enumerate}
        \item Les trois médiatrices d'un triangle se coupent au centre du cercle inscrit.
        \item La hauteur issue d'un sommet est perpendiculaire au côté opposé.
        \item Le centre de gravité est le point d'intersection des médianes.
        \item L'orthocentre est le point d'intersection des bissectrices.
    \end{enumerate}
    \item \textbf{Construction :} Soit $ABC$ un triangle tel que $AB=6cm$, $AC=8cm$ et $BC=7cm$.
    \begin{enumerate}
        \item Construire le triangle $ABC$ et tracer en \textbf{rouge} ses trois hauteurs. \textbf{(1,5 pt)}
        \item Nommer $H$ leur point d'intersection. Comment appelle-t-on ce point ? \textbf{(0,5 pt)}
        \item Tracer en vert le cercle circonscrit à ce triangle. \textbf{(1 pt)}
    \end{enumerate}
\end{enumerate}

\section*{\underline{Exercice 4 :} Droite des milieux (5 points)}
Soit $ABC$ un triangle. On appelle $I$ le milieu du segment $[AB]$ et $J$ le milieu du segment $[AC]$.
\begin{enumerate}
    \item Faire une figure complète. \textbf{(1 pt)}
    \item Que peut-on dire des droites $(IJ)$ et $(BC)$ ? Justifier par une propriété du cours. \textbf{(1,5 pt)}
    \item Si $BC = 9 cm$, calculer la longueur du segment $[IJ]$. \textbf{(1 pt)}
    \item Soit $K$ le symétrique de $I$ par rapport à $B$. La droite passant par $B$ et parallèle à $(AC)$ coupe $(JK)$ en $M$. (Question bonus : démontrer que $B$ est le milieu de $[IK]$). \textbf{(1,5 pt)}
\end{enumerate}

\end{document}